\section{Evaluation of authority and aligned behavior}\label{sec:C29:evaluation}
An authority evaluation needs tasks where the desired outcome and permitted action set are explicit. Include legitimate completion, conflicting instructions, expired grants, revoked access, ambiguous requests, delegated workers and pending effects. The evaluator should observe actual tool dispatch and environment changes, not only the final answer. A policy that verbally refuses but still causes an unauthorized effect has failed the operational constraint.

Separate model judgment from enforcement. Measure whether the model proposes appropriate actions, whether the runtime blocks disallowed ones, and whether useful work completes. This decomposition identifies complementary improvements. Better model behavior can reduce review and rejection load; stronger enforcement can contain residual mistakes. A combined score without these components cannot explain which layer supplies reliability.

Human-AI collaboration frameworks can clarify role expectations and degrees of agency~\cite{corpus250814825}. Their conceptual categories should not be treated as empirical proof of trust, shared intentionality or appropriate delegation. Evaluate the actual workflow: what the human understands, what the agent can change, and whether authority remains legible after interruptions and handoffs. Anthropomorphic language can otherwise conceal concrete responsibility gaps.

The statistical design should preserve user and organization clusters, because policies and preferences are shared within them. Repeated prompts against one policy do not establish cross-organization generalization. Include heterogeneous task families and report failures by authority type and consequence. Adaptive tests should search for boundary violations under declared budgets while benign tests measure overblocking and review burden.

The conclusion should be a bounded operational claim. A system can demonstrate useful behavior under specified objectives, constraints, principals and environments. It cannot infer unrestricted alignment from a high benchmark average, a constitutional prompt or a human-review stage. The scientific task is to make those conditions explicit, measure where they hold, and retain evidence about where they fail.
